\chapter{Literature Review}
\label{chap:literature_review}

The literature review maps the body of knowledge relevant to the research problem and
systematically identifies the gap that motivates this thesis. Following a structured search
strategy, the review synthesises seven thematic clusters: (A) CAPEX portfolio management and
strategic alignment, (B) project preemption and suspension mechanisms, (C) resource-constrained
project scheduling with skill considerations, (D) prescriptive analytics and AI-driven decision
support, (E) multi-criteria decision making for project prioritisation, (F) workforce and skill
demand forecasting, and (G) hybrid AI combining symbolic and subsymbolic methods. The chapter
closes with a cross-cluster gap analysis that directly motivates the thesis research questions.

% ─────────────────────────────────────────────────────────────────────────────
\section{Search Strategy}
\label{sec:search_strategy}

The systematic search was conducted on \textbf{Scopus} as the primary academic database,
supplemented by targeted searches on IEEE Xplore, ACM Digital Library, and INFORMS PubsOnline
for conference and journal papers not fully indexed at the time of search. The search applied
Boolean keyword combinations across seven thematic clusters, as summarised in
Table~\ref{tab:keywords}. A two-stage screening procedure was applied: title and abstract
screening followed by full-text relevance assessment. Only papers rated Medium or High
relevance against the research questions were retained for synthesis.

\begin{table}[htbp]
  \centering
  \caption{Keyword clusters used in the systematic Scopus search}
  \label{tab:keywords}
  \small
  \begin{tabular}{@{}lp{10.5cm}@{}}
    \toprule
    \textbf{Cluster} & \textbf{Search Terms} \\
    \midrule
    A — Portfolio & \texttt{CAPEX portfolio management}, \texttt{capital expenditure
      portfolio}, \texttt{project portfolio strategic alignment},
      \texttt{portfolio value optimisation} \\[0.3em]
    B — Preemption & \texttt{project preemption}, \texttt{preemptive RCPSP},
      \texttt{activity splitting}, \texttt{project suspension resumption} \\[0.3em]
    C — RCPSP & \texttt{resource-constrained project scheduling},
      \texttt{multi-skill scheduling}, \texttt{multi-project scheduling},
      \texttt{skill-based resource allocation} \\[0.3em]
    D — Prescriptive & \texttt{prescriptive analytics}, \texttt{decision support
      resource allocation}, \texttt{from predictive to prescriptive} \\[0.3em]
    E — MCDM & \texttt{MCDM project prioritisation}, \texttt{AHP TOPSIS portfolio},
      \texttt{dynamic weighting strategic alignment} \\[0.3em]
    F — Forecasting & \texttt{workforce skill demand forecasting},
      \texttt{temporal skill modelling}, \texttt{labour demand machine learning} \\[0.3em]
    G — Hybrid AI & \texttt{hybrid AI knowledge graph ontology},
      \texttt{neuro-symbolic AI decision support}, \texttt{ESCO skill ontology} \\
    \bottomrule
  \end{tabular}
\end{table}

% ─────────────────────────────────────────────────────────────────────────────
\section{CAPEX Portfolio Management and Strategic Alignment}
\label{sec:lit_capex}

Project portfolio management (PPM) has evolved from a purely financial discipline into a
strategic governance function concerned with selecting, prioritising, and terminating projects
in ways that maximise organisational value \citep{hansen2022seven}. In a comprehensive
bibliometric analysis spanning seven decades of PPM research, \citet{hansen2022seven} identify
resource allocation under competing constraints and the dynamic re-alignment of portfolios to
shifting strategic objectives as the two most persistently unresolved challenges in the field.
This observation directly contextualises the present research: in CAPEX-intensive environments,
resource conflicts among concurrent projects are structurally unavoidable, yet the literature
offers no established automated mechanism for resolving them in a strategically consistent
manner.

\citet{hansen2023principles} extend this line of inquiry by deriving seventeen normative PPM
principles organised around four semantic dimensions — portfolio value, strategic fit, balance,
and adaptive capacity. Their work provides a principled basis for evaluating automated
suspension-and-resumption recommendations: such recommendations should demonstrably improve
strategic fit and portfolio balance, not merely minimise makespan or resource cost.

At the level of individual portfolio optimisation models, \citet{chiang2013strategic} propose
an integer programming formulation for IT project portfolios that integrates a strategic
alignment score alongside expected benefits, development costs, and cross-project synergies.
Their model demonstrates that explicit strategic scoring, rather than implicit proxy measures
such as financial return, is essential for producing portfolios that genuinely reflect
organisational priorities. The same principle applies to CAPEX portfolios, where strategic
objectives — energy transition targets, regulatory compliance milestones, market positioning —
are often only loosely coupled with financial metrics. More recently, \citet{rodriguezgarcia2025ai}
extend this direction by combining machine learning with mathematical programming to
simultaneously optimise for strategic factors and synergies; their empirical results show that
strategy-aware portfolio selection consistently outperforms purely financial approaches, even
under uncertainty.

The growing complexity of CAPEX programmes — driven by regulatory pressures, talent scarcity,
and technological change — is reflected in the most recent contribution by
\citet{rodriguescoelho2025strategic}, who apply Structural Equation Modelling and Necessary
Condition Analysis to demonstrate that PPM is a necessary but not sufficient mediator between
organisational flexibility and value creation. The implication for AI-driven portfolio
governance is significant: strategic alignment must be embedded not only as an objective
function but as a hard constraint that any automated recommendation must respect.

% ─────────────────────────────────────────────────────────────────────────────
\section{Project Preemption and Suspension Mechanisms}
\label{sec:lit_preemption}

The concept of temporarily suspending an ongoing project activity to release a resource for
a higher-priority demand has a formal treatment in the operations research literature under
the label of \textit{preemption}. The foundational study by \citet{ballestin2008preemption}
demonstrates that allowing even a single preemptive interruption within a resource-constrained
project scheduling problem (RCPSP) yields statistically significant makespan reductions
relative to purely non-preemptive schedules. This result establishes preemption as an
operationally meaningful mechanism, not merely a theoretical relaxation.

\citet{vanpeteghem2010genetic} extend this analysis to multi-mode RCPSP, showing that a
genetic algorithm capable of selecting between preemptive and non-preemptive execution modes
consistently produces superior schedules when resource availability is uneven across the
planning horizon. Their work is particularly relevant to CAPEX portfolios, where the demand
profile for specialised engineering skills fluctuates significantly across project lifecycle
phases.

The most practically relevant model for the present thesis is provided by
\citet{hatamimoghaddam2024robust}, who introduce a robust multi-skill, multi-mode RCPSP
formulation with partial preemption. In their model, an activity may be suspended precisely
when a required skill set becomes unavailable within a specified time window, and resumed
when the skill pool is replenished. This mechanism maps directly onto the
suspension-and-resumption logic proposed in the thesis, with specialised engineering skill
shortages serving as the trigger condition.

\citet{vanhoucke2025impact} contribute a bi-criteria analysis examining the trade-off between
makespan reduction and the number of interruptions induced by preemption in problems with
time-varying resources. Their finding that optimal suspension frequency is sensitive to
resource availability patterns suggests that a data-driven, dynamic policy — as proposed by
the prescriptive AI framework in this thesis — is preferable to a static threshold-based rule.

Collectively, this body of work establishes that project suspension is a formally sound and
empirically effective mechanism for releasing constrained resources, yet existing models
operate on single projects or small portfolios and employ rule-based or combinatorial triggers
rather than strategic priority scoring. This is a first dimension of the research gap.

% ─────────────────────────────────────────────────────────────────────────────
\section{Resource-Constrained Project Scheduling with Skill Considerations}
\label{sec:lit_rcpsp}

The resource-constrained project scheduling problem (RCPSP) is one of the most extensively
studied problems in operations research. \citet{artigues2025fifty} trace fifty years of RCPSP
research, documenting the progression from exact branch-and-bound methods to metaheuristics,
and from single-project formulations to the multi-skill, multi-project variants most relevant
to contemporary CAPEX portfolio management. \citet{hartmann2022updated} provide a complementary
decade-level update, identifying skill-based and multi-project RCPSP as the fastest-growing
research sub-areas.

For portfolios of the scale addressed in this thesis (300+ concurrent projects),
\citet{gomez2023survey} demonstrate that shared resource pools across concurrent projects
are the primary source of scheduling conflicts and that the resulting \textit{resource-constrained
multi-project scheduling problem} (RCMPSP) is computationally intractable for exact methods
at real-world scale. This motivates the use of heuristic and data-driven approaches.

The modelling of human skill heterogeneity within RCPSP has received increasing attention.
\citet{snauwaert2021algorithm} introduce a distinction between \textit{skill breadth} (the
number of distinct skills an individual possesses) and \textit{skill depth} (proficiency level
per skill), demonstrating that both dimensions independently affect assignment quality in a
genetic algorithm framework. \citet{snauwaert2023classification} formalise a classification
scheme for multi-skilled RCPSP (MSRCPSP) and release standardised benchmark instances, providing
the empirical infrastructure needed for rigorous model evaluation. Their most recent contribution
\citep{snauwaert2025hierarchical} introduces six MSRCPSP variants with hierarchical skill
structures — where junior skills are subsumed by senior ones — directly mirroring the tiered
engineering competency pools typical of large CAPEX programmes.

The industrial feasibility of solving large-scale MSRCMPSP instances is demonstrated by
\citet{torba2024reallife}, who address a real-world scheduling problem at SNCF heavy
maintenance involving thousands of activities, multiple skill types, and weighted tardiness
objectives. Their memetic algorithm solution approach is the closest real-world analogue to
the computational context of the present thesis. Similarly, \citet{haroune2023multiproject}
address multi-project assignment under shared employee skill pools with efficiency levels,
minimising both weighted tardiness and goal deviations — a formulation that naturally
accommodates project priority weights, linking Cluster C to the MCDM considerations of
Cluster E.

% ─────────────────────────────────────────────────────────────────────────────
\section{Prescriptive Analytics and AI-Driven Decision Support}
\label{sec:lit_prescriptive}

Prescriptive analytics denotes the class of analytical methods that go beyond description
and prediction to generate actionable recommendations — that is, explicit decisions or
interventions designed to achieve a specified objective under identified constraints
\citep{lepenioti2020prescriptive}. In their systematic review of 56 key papers,
\citet{lepenioti2020prescriptive} identify resource allocation as the primary application
domain of prescriptive methods and distinguish three core technical approaches: optimisation,
simulation, and heuristic rule learning. \citet{frazzetto2019prescriptive} complement this
picture from a data management perspective, identifying constraint satisfaction, optimisation,
and simulation as the architectural pillars of any prescriptive analytics system — a
decomposition that maps directly onto the design components of the proposed framework.

The theoretical foundation for embedding prediction within optimisation is provided by
\citet{bertsimas2020predictive}, who formalise the transition from predictive to prescriptive
analytics and introduce the \textit{coefficient of prescriptiveness} as a performance measure
capturing the value added by incorporating predictions into decisions. Their framework is
particularly relevant to the thesis's analytics chain: temporal skill demand forecasts (the
predictive layer) are the principal input to the suspension-and-resumption optimisation (the
prescriptive layer).

\citet{wissuchek2025prescriptive} offer the most recent and comprehensive systems-theoretic
treatment of prescriptive analytics systems (PAS), identifying 23 architectural components
including data pipeline, intervention model, feedback mechanism, and explainability layer.
Their analysis reveals that few existing PAS implementations address multi-objective portfolio
decisions with dynamic constraint sets — precisely the configuration required for CAPEX
resource governance. \citet{misic2020analytics} further contextualise prescriptive methods
within the broader operations management literature, noting that their comparative advantage
over predictive methods is most pronounced when decision windows are short and resource
constraints are binding, which is consistently the case in CAPEX portfolio management.

% ─────────────────────────────────────────────────────────────────────────────
\section{Multi-Criteria Decision Making for Project Prioritisation}
\label{sec:lit_mcdm}

Project prioritisation in the context of competing resource constraints is inherently a
multi-criteria problem: strategic importance, financial return, regulatory urgency,
stakeholder commitment, and execution risk are incommensurable attributes that must be
integrated into a single ranking or weighting scheme. \citet{kandakoglu2024mcdm} conduct a
systematic literature review of MCDM methods applied to project portfolio selection, covering
the Analytic Hierarchy Process (AHP), PROMETHEE, TOPSIS, and Data Envelopment Analysis (DEA).
A central finding of their review is that \textit{dynamic re-prioritisation under constraint
changes} is an open research problem: existing MCDM applications to PPM are almost uniformly
static, producing a priority ordering at a single point in time rather than continuously
updating weights as resource availability and strategic context evolve. This constitutes a
second dimension of the research gap.

\citet{chatterjee2018fuzzy} address the treatment of linguistic uncertainty in expert
prioritisation judgements through fuzzy AHP. Their contribution is methodologically important
for the present thesis: criteria weights assigned by portfolio managers are inherently
imprecise and may change over time, requiring a weighting framework that accommodates
vagueness without reducing to arbitrary scores. The dynamic extension of this approach — where
fuzzy weights are recalibrated in response to changing resource conditions and project progress
— forms the conceptual basis of Sub-Question 2.

The application of MCDM to capital-intensive industrial portfolios is examined by
\citet{alhashimi2017mcdm}, who compare the output coherence of AHP, PROMETHEE, and TOPSIS
applied to oilfield CAPEX portfolio prioritisation. Their key finding — that method output
coherence is high when criteria weights are well-calibrated but diverges under expert
disagreement — underscores the need for systematic weight elicitation and sensitivity
analysis in any automated de-prioritisation system.

% ─────────────────────────────────────────────────────────────────────────────
\section{Workforce and Skill Demand Forecasting}
\label{sec:lit_forecasting}

Accurately modelling the temporal demand for specialised skill pools is a prerequisite for
any anticipatory suspension-and-resumption mechanism. Without a reliable forecast of when
and where skill shortages will materialise, suspension recommendations can only be generated
reactively — after conflicts have already disrupted project execution. \citet{safarishahrbijari2018workforce}
provide a systematic review of workforce forecasting models, covering time series methods,
Markov chain approaches, system dynamics, and discrete event simulation. Their taxonomy
reveals a consistent pattern: models that integrate demand-side dynamics (project workload
profiles) with supply-side constraints (skill availability, attrition, hiring lead times)
outperform demand-only models in medium-to-long forecast horizons. For a 300+ project
portfolio with overlapping lifecycle phases, such integrated modelling is essential.

At the intersection of deep learning and skill demand forecasting, \citet{macedo2022skills}
apply recurrent neural networks (LSTM) to a decade of monthly skill demand observations,
demonstrating that multivariate models capturing inter-skill correlations outperform
univariate baselines in multi-step forecasting tasks. Their work provides empirical evidence
that the temporal skill demand patterns required by Sub-Question 1 are learnable from
historical project staffing data, subject to sufficient data quality and volume.

Temporal Fusion Transformers \citep{lim2021tft} represent the current state of the art for
multi-horizon time series forecasting with mixed-frequency inputs and static covariates.
Their architecture — combining multi-head attention for long-range dependencies with
interpretable variable selection — is a natural candidate for skill demand forecasting in
CAPEX portfolios, where project metadata (type, phase, contractual milestones) constitutes
structured static covariates that modulate temporal demand patterns.

% ─────────────────────────────────────────────────────────────────────────────
\section{Hybrid AI: Symbolic, Subsymbolic, and Their Integration}
\label{sec:lit_hybrid_ai}

A recurring limitation of purely data-driven approaches in enterprise decision support is
their dependence on the statistical regularity of historical data. In domains characterised
by low data volume, high structural complexity, or rapidly changing operational contexts —
all of which apply to CAPEX portfolio management — subsymbolic models alone are insufficient.
The paradigm of \textit{hybrid AI} addresses this limitation by combining the complementary
strengths of symbolic and subsymbolic methods \citep{garcez2022neurosymbolic}.

\textbf{Symbolic AI} encompasses methods that operate on explicit, human-interpretable
representations of knowledge: ontologies, rule systems, logic programs, and knowledge graphs.
Symbolic methods are inherently interpretable, capable of encoding domain constraints
formally, and effective in low-data regimes. Their weakness is brittleness: they cannot
generalise beyond explicitly encoded knowledge and fail gracefully when assumptions are
violated.

\textbf{Subsymbolic AI} encompasses statistical learning methods — including classical machine
learning, deep learning, and generative AI — that extract patterns from data without explicit
symbolic encoding. Subsymbolic models are powerful pattern recognisers that generalise to
unseen inputs, but they lack interpretability, require large training sets, and cannot natively
incorporate hard domain constraints.

\textbf{Hybrid AI} combines both paradigms, typically by using structured symbolic knowledge
to constrain, enrich, or initialise subsymbolic learning processes. Gartner identifies hybrid
AI as an emerging high-value technology trend in enterprise AI adoption, highlighting its
particular relevance for domains where domain expertise must be preserved while adapting to
changing data distributions \citep{gartner2024hype}. \citet{garcez2022neurosymbolic} formalise
this integration as \textit{neurosymbolic AI}, demonstrating that combining neural learning
with symbolic reasoning yields systems that are simultaneously more accurate and more
interpretable than either component in isolation.

In the context of this thesis, the hybrid AI architecture serves three distinct roles. First,
a symbolic layer encodes domain knowledge about skill taxonomies, project types, and
organisational constraints in a structured, machine-processable form. Second, a subsymbolic
layer learns temporal demand patterns from historical project staffing data, with the symbolic
layer providing structured features and constraint boundaries. Third, a prescriptive layer
integrates both outputs into suspension-and-resumption recommendations that are both
data-driven and domain-knowledge-compliant.

\textbf{The ESCO Ontology} provides the symbolic backbone for skill representation in this
framework. Developed by the European Commission, ESCO (European Skills, Competences,
Qualifications and Occupations) is a multilingual, hierarchical taxonomy of over 13{,}900
skills and competences formally linked to occupations and qualifications through a structured
ontological hierarchy \citep{le2014esco}. ESCO's formal structure — skills organised into
pillar hierarchies with defined relationships (\texttt{skos:broader},
\texttt{skos:narrower}, \texttt{skos:related}) — makes it directly suitable as a semantic
backbone for skill-demand modelling in CAPEX portfolios. Specifically, ESCO enables consistent
skill labelling across heterogeneous project data sources, ontological inference of skill
substitutability (i.e., which alternative skills can fulfil a demand when the primary skill
is unavailable), and semantic enrichment of subsymbolic demand forecasting models through
ontology embeddings.

The AI-4-SME Framework \citep{peter2025ai4sme}, developed at FHNW School of Business,
provides the overarching methodological structure for translating the hybrid AI architecture
into an implementable solution. The framework's Design--Build--Run methodology distinguishes
Knowledge-Intensive Activities (KIAs) — which require contextual judgement, problem solving,
and expert knowledge — from Data-Intensive Activities (DIAs) — which require high-volume
data processing and statistical learning. In the CAPEX resource management context, demand
forecasting and strategic priority scoring are KIAs; allocation optimisation and constraint
propagation are DIAs. The framework's hybrid build path, combining knowledge-based systems
(Neo4j, GraphDB, Stardog) with data-based ML pipelines (TensorFlow, PyTorch), maps directly
onto the proposed framework architecture.

% ─────────────────────────────────────────────────────────────────────────────
\section{Research Gap and Positioning}
\label{sec:lit_gap}

The synthesis of the seven clusters reveals a coherent set of gaps at the intersection of the
reviewed domains. Table~\ref{tab:gap} summarises the state of knowledge per dimension and
identifies the contribution of the present thesis.

\begin{table}[htbp]
  \centering
  \caption{Research gap summary across literature clusters}
  \label{tab:gap}
  \small
  \begin{tabular}{@{}p{4.2cm}p{5.5cm}p{4.5cm}@{}}
    \toprule
    \textbf{Dimension} & \textbf{State of the Art} & \textbf{Gap / Thesis Contribution} \\
    \midrule
    Strategic alignment in PPM &
    Static portfolio selection models with strategic scoring
    \citep{chiang2013strategic, rodriguezgarcia2025ai} &
    No dynamic, AI-driven re-alignment triggered by real-time skill shortages \\[0.4em]

    Preemptive project scheduling &
    Rule-based and combinatorial preemption models for single or small portfolios
    \citep{ballestin2008preemption, hatamimoghaddam2024robust} &
    No prescriptive AI layer translating portfolio-level strategic priorities into
    suspension triggers \\[0.4em]

    Multi-skill RCPSP at scale &
    Heuristic and metaheuristic methods for industrial-scale instances
    \citep{torba2024reallife, snauwaert2025hierarchical} &
    No integration with strategic MCDM weighting or AI-driven priority scoring \\[0.4em]

    Prescriptive analytics &
    Theoretical frameworks and isolated domain applications
    \citep{lepenioti2020prescriptive, wissuchek2025prescriptive} &
    No PAS implementation for multi-project CAPEX suspension-and-resumption
    decisions \\[0.4em]

    MCDM for project prioritisation &
    Static, point-in-time ranking with expert-elicited weights
    \citep{kandakoglu2024mcdm, chatterjee2018fuzzy} &
    No dynamic, continuously updated MCDM weighting framework responsive to
    resource availability and project progress \\[0.4em]

    Skill demand forecasting &
    Temporal models for labour market or HR demand
    \citep{safarishahrbijari2018workforce, macedo2022skills} &
    No application to project-level, portfolio-wide skill pool demand across
    300+ concurrent CAPEX projects \\[0.4em]

    Hybrid AI in enterprise decision support &
    Neurosymbolic frameworks and ESCO-based skill extraction
    \citep{garcez2022neurosymbolic, gartner2024hype} &
    No hybrid AI system combining ESCO ontology, temporal demand forecasting,
    and prescriptive suspension optimisation in a unified CAPEX framework \\
    \bottomrule
  \end{tabular}
\end{table}

The overarching research gap can be stated as follows: \textit{no existing work unifies
(1)~strategic priority scoring via dynamic MCDM, (2)~AI-driven temporal modelling of
specialised skill pool demand, and (3)~prescriptive suspension-and-resumption optimisation
into a single, operationally deployable framework for CAPEX portfolio management.} Each of
the three components exists in isolation within its respective literature cluster; their
integration, guided by the AI-4-SME framework and grounded in ESCO-based hybrid AI
architecture, constitutes the novel contribution of this thesis.
